\documentclass[11pt,a4paper]{article}
\usepackage{../pelm}
\begin{document}
\paper{Quiz 3}{Weeks 5--6: Style, Tone and Format; Building Without Code}{10}{10 minutes}

\begin{questions}

\q{Why does the instruction ``be professional'' produce inconsistent results?}
  {The model does not understand the word}
  {It describes a feeling rather than a text, and is compatible with many different outputs}
  {It is too short to affect the output}
  {Professional writing is rare in training data}

\q{In the three dimensions of register, \textbf{tenor} refers to:}
  {The subject matter and its vocabulary}
  {The relationship between writer and reader}
  {Whether the text is written or spoken}
  {The emotional warmth of the text}

\q{Which of these is a \textbf{behaviour} rather than an adjective, and therefore checkable?}
  {Make it warm}
  {Make it reassuring}
  {Use the reader's name and acknowledge their situation}
  {Make it feel more human}

\q{You ask for a table with a ``deadline'' column, but the source document states no deadline.
   Without a fallback instruction, what is most likely?}
  {The cell is left empty}
  {The system reports an error}
  {A plausible deadline appears in the cell}
  {The whole table is refused}

\q{Why does adding ``if a field cannot be filled from the source, write NOT STATED'' work?}
  {It appeals to the model's honesty}
  {It makes the honest answer a likely continuation}
  {It disables the model's creativity}
  {It reduces the temperature}

\q{Which is a well-formed negative constraint?}
  {Do not be boring}
  {Do not write badly}
  {Do not use any sentence longer than 20 words}
  {Do not disappoint the reader}

\q{When generating a working browser tool from a description, what is actually happening?}
  {The system assembles pre-built components from a library}
  {The system predicts code, in the same way it predicts prose}
  {The system searches for an existing tool and copies it}
  {The system compiles your description into a program}

\q{Why do generated tools succeed most reliably on ordinary tasks?}
  {Ordinary tasks require less code}
  {Common tasks appear thousands of times in near-identical form in training data}
  {The model tests ordinary tasks before returning them}
  {Ordinary tasks have fewer edge cases by definition}

\q{A generated tool has a clean interface, working buttons, and produces a confidently formatted
   number that is wrong every time. Which failure is this?}
  {Silently discarded input}
  {Correct only for the cases tried}
  {Wrong logic with right appearance}
  {Context overflow}

\q{Which statement best captures the central risk of building a tool you cannot read?}
  {It may be slow}
  {You can judge whether it looks right, but not whether it is right}
  {It may stop working when the service changes}
  {It cannot handle large amounts of data}

\end{questions}
\end{document}
